\documentclass[11pt, a4paper]{article}

% --- PACCHETTI ---
\usepackage[utf8]{inputenc}
\usepackage[italian]{babel}
\usepackage[margin=2.5cm]{geometry} % Imposta i margini della pagina
\usepackage{graphicx}               % Per inserire il logo
\usepackage{tabularx}               % Per tabelle larghe quanto la pagina
\usepackage[table]{xcolor}          % Per colorare le celle delle tabelle
\usepackage[hidelinks]{hyperref}    % Per stile link cliccabili
\usepackage{lmodern}      % Rende i font scalabili all'infinito
\usepackage{anyfontsize}  % Permette di usare dimensioni personalizzate

\begin{document}

% ==============================================================================
% FRONTESPIZIO (PAGINA 1)
% ==============================================================================
\begin{titlepage}
    \centering
    \vspace*{0.5cm}

    % LOGHI
    \includegraphics[width=0.45\textwidth]{assets/logo_unipd.jpg}
    \hfill
    \includegraphics[width=0.45\textwidth]{assets/logo.png}\par
    \vspace{0.5cm}

    % NOME DEL GRUPPO
    {\fontsize{45}{50}\sffamily\selectfont \textbf{On-By-One}\par} 
    \vspace{1cm}

    % TITOLO DEL DOCUMENTO
    {\huge \bfseries Piano di progetto \par}
    
    \vspace{1.5cm}

    % MEMBRI DEL GRUPPO
    {\large
    \textbf{On-By-One}:\par
    \vspace{0.3cm}
    \begin{tabular}{l l}
        Sara Ballardin & 2101048 \\
        Shuhui Chen & 2138001 \\
        Simone Giuliano & 2145591 \\
        Alvise Greggio & 2141004 \\
        Daniele Nuzzaci & 2139039 \\
        Federico Temporin & 2146764 \\
        Tommaso Zorzetto & 2148629 \\
    \end{tabular}
    \par}

    \vfill

    % TABELLA INFORMAZIONI DOCUMENTO E DESTINATARI
    \begin{table}[h]
        \centering
        \arrayrulecolor{gray!50}
        \renewcommand{\arraystretch}{1.5}
        \begin{tabularx}{\textwidth}{|c|c|c|X|}
            \hline
            \rowcolor{gray!10} \multicolumn{4}{|c|}{\textbf{Informazioni documento}} \\
            \hline
            \rowcolor{gray!10} \textbf{Versione} & \textbf{Data} & \textbf{Stato} & \multicolumn{1}{c|}{\textbf{Destinatari}} \\
            \hline
            - & - & Bozza & 
            On-By-One \newline 
            Prof. Tullio Vardanega \newline 
            Prof. Riccardo Cardin \newline
            Vimar S.P.A. \\
            \hline
        \end{tabularx}
        \arrayrulecolor{black}
    \end{table}
    
    \vspace{0.5cm}
    
    % CONTATTI
    {\large Contatti: \href{mailto:onbyone.swe@gmail.com}{onbyone.swe@gmail.com} }
    
\end{titlepage}

\newpage

% ==============================================================================
% REGISTRO DELLE MODIFICHE (PAGINA 2)
% ==============================================================================
\section*{Registro delle modifiche}

\vspace{1cm}
\begin{center}
    \arrayrulecolor{gray!50}
    \renewcommand{\arraystretch}{1.5}
    \begin{tabularx}{\textwidth}{|c|c|c|X|c|}
        \hline
        \rowcolor{gray!10} 
        \multicolumn{1}{|c|}{\textbf{Versione}} & 
        \multicolumn{1}{c|}{\textbf{Data}} & 
        \multicolumn{1}{c|}{\textbf{Autore}} & 
        \multicolumn{1}{c|}{\textbf{Descrizione}} & 
        \multicolumn{1}{c|}{\textbf{Verificatore}} \\
        \hline
        - & - & - & - & - \\
        \hline
        0.1.0 & 2026/11/09 & Sara Ballardin & Definizione struttura del documento & - \\
        \hline
    \end{tabularx}
    \arrayrulecolor{black}
\end{center}
\vfill
\newpage

% ==============================================================================
% INDICE (PAGINA 3)
% ==============================================================================
\tableofcontents
\newpage


% ==============================================================================
% 1. INTRODUZIONE
% ==============================================================================
\section{Introduzione}

\subsection{Scopo del documento}
Descrivere la finalità del Piano di Progetto: pianificare le attività operative, definire la scansione temporale, ripartire le risorse orarie ed economiche tra i ruoli e monitorare l'avanzamento dei lavori per tutta la durata del progetto.

\subsection{Scopo del prodotto}
Riassumere in breve gli obiettivi del sistema da realizzare, evidenziando il problema che risolve, i destinatari e il valore atteso per la proponente.

\subsection{Glossario}
Spiegare la convenzione tipografica adottata per evidenziare i termini tecnici, specialistici o ambigui e rimandare al documento esterno di Glossario.

\subsection{Riferimenti}
\subsubsection{Riferimenti normativi}
Elencare i riferimenti vincolanti: regolamento del progetto didattico, Norme di Progetto, capitolato d'appalto e documento di candidatura (dichiarazione degli impegni).

\subsubsection{Riferimenti informativi}
Elencare le fonti consultate per la gestione del progetto: slide del corso, testi su metodologie agili e documentazione tecnica di supporto.

\newpage


% ==============================================================================
% 2. ANALISI DEI RISCHI
% ==============================================================================
\section{Analisi dei rischi}
Descrivere la procedura con cui vengono identificati, classificati e monitorati i rischi (codice identificativo, nome, descrizione, misure di mitigazione, probabilità di occorrenza e livello di impatto).

\subsection{Rischi tecnologici}
Descrivere i rischi legati all'apprendimento e all'integrazione di nuove tecnologie, alla comparsa di anomalie nel codice e a procedure di collaudo non sufficienti.

\subsection{Rischi interpersonali}
Descrivere i rischi legati a dinamiche relazionali interne, opinioni contrastanti sul lavoro da svolgere e disallineamenti di motivazione o partecipazione.

\subsection{Rischi organizzativi}
Descrivere i rischi legati a un'errata distribuzione del carico di lavoro, alla stima imprecisa delle ore per ruolo (sovrastima o sottostima) e al mancato rispetto delle scadenze interne o verso gli stakeholder.

\subsection{Rischi legati a requisiti e comunicazione}
Descrivere i rischi derivanti da una comprensione incompleta o ambigua dei requisiti del capitolato e da un mancato allineamento con la proponente o con i committenti.

\subsection{Rischi personali}
Descrivere i rischi derivanti da impegni accademici (sessioni d'esame) o imprevisti personali dei singoli membri che possono ridurre la disponibilità oraria nel periodo.

\newpage

% ==============================================================================
% 3. PIANIFICAZIONE GENERALE
% ==============================================================================
\section{Pianificazione generale}

\subsection{Modello di sviluppo}
Descrivere il ciclo di vita del software adottato e il modello iterativo a periodi di avanzamento (sprint), illustrando il ciclo di pianificazione, esecuzione, verifica, consuntivazione e retrospettiva.

\subsection{Gestione e rotazione dei ruoli}
Descrivere brevemente le responsabilità assegnate a ciascun ruolo, i costi orari fissati da regolamento e lo schema di rotazione periodica adottato per assicurare equità di impegno e assenza di conflitti di interesse durante le verifiche.

\subsection{Scadenziario di massima}
Riportare le scadenze cardine del progetto: traguardi intermedi verso la Requirements and Technology Baseline (RTB), verso la Product Baseline (PB) e data limite di consegna finale, accompagnate da un diagramma di Gantt (o PERT) complessivo.

\subsection{Pianificazione economica iniziale}
Riportare il prospetto generale definito in fase di candidatura: totale ore produttive per persona e per ruolo, ripartizione del monte ore nei macro-periodi e preventivo di spesa complessivo vincolante.

\subsection{Meccanismi di controllo e rendicontazione}
Descrivere le modalità con cui il gruppo monitora a ogni fine sprint l'avanzamento effettivo e le risorse spese rispetto al preventivo iniziale, aggiornando il consuntivo di periodo e il preventivo a finire.

\newpage

% ==============================================================================
% 4. PERIODI DI AVANZAMENTO
% ==============================================================================
\section{Periodi di avanzamento}
Introdurre la sezione incrementale destinata a documentare, per ogni singolo ciclo di avanzamento, le attività programmate, i risultati conseguiti, i rischi manifestati, il consuntivo orario ed economico e le decisioni di retrospettiva.

\subsection{Periodo di candidatura}
Rendicontare sinteticamente le attività preparatorie svolte prima dell'aggiudicazione dell'appalto: formazione del team, canali di comunicazione, analisi dei capitolati, redazione della candidatura e primi rischi affrontati.

% ------------------------------------------------------------------------------
% STRUTTURA TIPO PER OGNI SPRINT
% ------------------------------------------------------------------------------
\subsection{Sprint N [Data inizio -- Data fine]}

\subsubsection{Attività programmate}
Elencare le attività pianificate per il periodo, specificando gli obiettivi operativi assegnati ai membri nei rispettivi ruoli.

\subsubsection{Attività svolte}
Descrivere le attività effettivamente completate nel corso del periodo, segnalando eventuali compiti non conclusi o posticipati.

\subsubsection{Rischi attesi e manifestati}
Confrontare i rischi previsti all'avvio del periodo con quelli realmente occorsi, illustrando le misure di mitigazione applicate e la loro efficacia.

\subsubsection{Consuntivo di periodo}
Riportare le tabelle di rendicontazione oraria ed economica del periodo: ore lavorate per membro e per ruolo, costi effettivi sostenuti e scostamenti rispetto a quanto preventivato.

\subsubsection{Preventivo a finire}
Presentare la tabella aggiornata con le ore e il budget residui per ciascun ruolo, verificando che la proiezione della spesa finale non superi il limite massimo prefissato.

\subsubsection{Retrospettiva}
Esporre le considerazioni emerse durante la riunione di fine periodo: punti di forza, difficoltà riscontrate nei processi o nelle stime e azioni correttive da applicare nel periodo successivo.

\end{document}